\section{Human Evaluation Procedure}~\label{appx:user}

This appendix describes the protocol for the human evaluation in Sec.~\ref{exp:human_eval}. It summarizes the participant instructions, trial procedure, and interview guide. The physical platforms and task domains match the VLM evaluation.

% Sources in asset/BRL_interview/docs
% Participant instructions from 참가자_실험매뉴얼.md.
% Trial setup and recording from 연구자_실험매뉴얼_v2.md.
% Core questions and optional prompts from 인터뷰.md.
% The protocol describes the scheduled study rather than verified completion.
\subsection{Participant Briefing}~\label{appx:user_briefing}
The briefing explains the purpose of the study, the sequence of activities, and the participation compensation. Participants complete the research participation and personal information consent forms before the tasks begin. They are informed that participation is voluntary and that they can request a break or withdraw during the session.

Participants receive explanations of the task goals, object categories, available robot actions, and tablet controls for both domains. They then watch example videos of task execution. The instructions ask participants to monitor the camera views and answer robot questions according to their own judgment. Participants are told that an incorrect answer does not interrupt the experiment and that they should select the answer they consider most appropriate if they are unsure.

\subsection{Domain Information}~\label{appx:user_scenarios}
In WasteSorting, participants receive information about the four waste categories and the corresponding bins. The available actions detect waste objects, pick an object, and place the held object in a selected bin. In TomatoHarvesting, participants receive information about ripe and unripe tomatoes and the distinction between fresh and rotten tomatoes. The available actions navigate between locations, detect tomatoes at a stem, pick a tomato, scan a held tomato, and place or discard the tomato. The task goal requires fresh ripe tomatoes to be harvested, rotten tomatoes to be discarded, and unripe tomatoes to remain on the stem.

Each WasteSorting trial contains four objects. Participants are taught to distinguish general waste, paper, plastic, and cans and to place each object in the corresponding bin. The interface marks the categories with blue, green, yellow, and gray boxes, respectively. Object identifiers follow detection order and a position rule if several objects are detected together. Participants are instructed that the identifiers do not indicate object categories.

Each TomatoHarvesting trial contains four tomatoes distributed across two stems. The briefing distinguishes three appearances corresponding to fresh ripe tomatoes, rotten ripe tomatoes, and fresh unripe tomatoes. Participants are taught that ripeness can be judged from a distance, whereas a closer inspection is needed to judge rot. The interface marks fresh ripe tomatoes in green, unripe tomatoes in yellow, and rotten tomatoes in red. The robot can detect and pick tomatoes only at its current stem. Tomato identifiers follow detection order, with the rightmost tomato assigned first if several tomatoes are detected together. The task ends after the tomatoes have been processed according to the task rules and the robot has returned to the charging station.

\subsection{Tablet Interaction}~\label{appx:user_assignment}
The tablet interface displays robot camera views and the current query. Ours presents a question about a selected state fact with Boolean response options. KnowNo presents candidate actions and requests an action choice. Participants choose a response through the tablet, and the robot continues execution after the response is submitted. Ours can request several state confirmations before the confidence threshold is met. KnowNo requests one action choice per trigger.

Participants start each task through the corresponding start button and use the response buttons on the right side of the screen to answer queries. Voice, gesture, and manual control functions are not used in the experiment. If a communication problem prevents interaction, execution is paused until the connection is restored.

\subsection{Trial Schedule}~\label{appx:user_schedule}
% Assigned sequences supplied on 2026-10-09 for the evaluation scheduled for October 12--15.
% Confirm the executed sequences and the participant ID mapping after data collection.
The planned evaluation comprises five participants with two trials per domain and method for each participant. This schedule yields $10$ trials per domain and method and $40$ trials in total. Table~\ref{tab:human_assignment} gives the assigned domain and method orders. Three participants start with WasteSorting and two with TomatoHarvesting. In each domain, three participants start with Ours and two with KnowNo.

Each participant completes four consecutive trials in one domain before the four trials in the other domain. The assigned sequences vary the domain and method orders across participants. With five participants, the orders are not fully balanced.

\subsection{Trial Procedure}~\label{appx:user_procedure}
A facilitator explains the tasks and manages the session, while an assistant prepares the robot and object arrangement for each assigned scenario. Before each trial, the experimenters reset the environment, check the selected method and scenario, and confirm that the camera views are available. The facilitator starts the tablet screen recording, and the experimenters check that the trial logs are saved after execution.

After each trial, the participant completes a questionnaire about the robot system just experienced. The questionnaire contains nine items about system understanding and performance, with responses on a seven point scale from strong disagreement to strong agreement. Participants are asked to evaluate the most recent trial rather than combine experiences from earlier trials. They can rest while the assistant prepares the next scenario. This sequence is repeated for all eight trials, followed by an interview about the complete experience.

At the end of the session, the experimenters save the tablet and external camera recordings and verify that the logs match the participant identifier, trial order, domain, and method. These records document task execution and the responses supplied through the interface.

\subsection{Interview Guide}~\label{appx:user_interview}
The interview guide invites participants to describe both favorable and unfavorable experiences with the two methods. Participants can report no difference or state that they do not remember a particular event. The interviewer asks all six core questions and selects additional prompts only if clarification or a concrete example is needed. The questions below are English translations of the Korean guide.

\begin{enumerate}
\item What was the most memorable difference between the two methods? Please also say if you noticed no difference.
\item For each method, what did you look at and think about before you chose an answer to a robot question? Please describe a question that you remember.
\item Did the information on the screen or the response options help you choose an answer or make it difficult?
\item How did you feel about the timing of the robot questions? Please describe a scene that you remember.
\item Across the entire task, how did the two methods differ in the effort you spent and your sense of task progress?
\item Which method would you choose for a similar task, and why? You may have no preference or choose differently depending on the situation.
\end{enumerate}

Additional prompts ask what participants checked before an answer, whether they considered task rules or subsequent robot actions, and whether they hesitated. Prompts about the interface ask whether the target object was clear, whether the camera views provided enough information, and whether the available options included the intended answer. Prompts about timing ask about opportunities to intervene, repeated questions, and interruptions to observation or thought. Prompts about effort distinguish the difficulty of an individual judgment from the effort of repeated responses and continued monitoring. The final prompts ask whether preferences differ across domains and which improvement participants would prioritize.

If further explanation is needed, the guide permits one or two common examples of questions from the two methods. The interviewer explicitly identifies these as prepared examples that may differ from the actual task experience. Participants describe the information and judgments needed to answer each example and indicate whether they recall a similar event during the experiment. The interviewer ends with a summary of the main responses and invites corrections or further comments.

% Replace the current human aggregates after the five-participant study.
% Confirm the executed procedure and final participant demographics after collection.
% The participant manual estimates 25--35 minutes for the interview, whereas the
% standalone guide gives a shorter question schedule. Report measured durations.
% Add the actual interview recording, transcription, and analysis procedure after
% confirmation. The analysis recommendations in the TOAST overview do not establish
% that a particular coding or analysis method was performed.
